\section{图着色与冲突安排}
\label{sec:04-Discrete-Mathematics/03-Graph-Theory/05}

编译器把一段函数编译成机器码时要决定：哪个变量放进哪个寄存器。寄存器只有十几个，函数里的变量可能有几十个，所以必须共用。共用的规矩只有一条——两个变量如果在某一刻同时``活着''（都还会被后面的代码读到），它们就不能占同一个寄存器，否则一个会把另一个覆盖掉。

上一节的做法在这里派得上用场：把每个变量画成顶点，把``活跃区间相交''画成边。可这一次两个顶点之间没有``跨组''可言，任何两个变量都可能冲突。要问的也不是配对，而是最少用几个寄存器能把所有变量安顿下来。

\subsection{冲突变成边，资源变成颜色}

给每个顶点涂一种颜色，要求相邻顶点异色。一种颜色就是一份资源——一个寄存器、一个考试时段、一段频谱。最少需要的颜色数记作 \(\chi(G)\)，叫色数。

\begin{center}
\begin{tikzpicture}[x=1cm,y=1cm,font=\small,>=stealth]
  \draw[black!30] (-0.05,0.45) -- (4.75,0.45);
  \node[black!55,font=\footnotesize] at (4.9,0.2) {时间};
  \draw[black!50,fill=blue!30] (0,2.44) rectangle (1.8,2.76);
  \draw[black!50,fill=orange!40] (0.9,1.84) rectangle (2.7,2.16);
  \draw[black!50,fill=blue!30] (2.25,1.24) rectangle (4.05,1.56);
  \draw[black!50,fill=green!30] (1.35,0.64) rectangle (4.5,0.96);
  \node[left] at (-0.15,2.6) {\(a\)};
  \node[left] at (-0.15,2.0) {\(b\)};
  \node[left] at (-0.15,1.4) {\(c\)};
  \node[left] at (-0.15,0.8) {\(d\)};
  \node[black!70,font=\footnotesize] at (2.2,3.35) {活跃区间};
  % 冲突图
  \draw[black!60] (6.3,2.6) -- (7.9,3.2);
  \draw[black!60] (6.3,2.6) -- (7.9,1.3);
  \draw[black!60] (7.9,3.2) -- (9.5,2.6);
  \draw[black!60] (7.9,3.2) -- (7.9,1.3);
  \draw[black!60] (9.5,2.6) -- (7.9,1.3);
  \draw[black!60,fill=blue!30] (6.3,2.6) circle (8pt);  \node at (6.3,2.6) {\(a\)};
  \draw[black!60,fill=orange!40] (7.9,3.2) circle (8pt);\node at (7.9,3.2) {\(b\)};
  \draw[black!60,fill=blue!30] (9.5,2.6) circle (8pt);  \node at (9.5,2.6) {\(c\)};
  \draw[black!60,fill=green!30] (7.9,1.3) circle (8pt); \node at (7.9,1.3) {\(d\)};
  \node[black!70,font=\footnotesize] at (7.9,3.9) {冲突图};
\end{tikzpicture}

\emph{四个变量的活跃区间画在左边，区间相交就在右边连一条边。\(b,c,d\) 两两相交，围成一个三角形，三者必须各占一个寄存器。\(a\) 与 \(c\) 的区间错开了，两人同色，于是三个寄存器就够——上界与下界在 3 这个数上碰头。}
\end{center}

颜色叫什么无关紧要，红蓝绿也好，寄存器编号也好，换个标签不改变任何结构。真正要回答的是两个方向的问题：能用多少种颜色涂完（上界，靠构造），以及少于多少种一定涂不完（下界，靠找障碍）。这两件事的思路截然不同，图上那句``三角形''和那句``\(a\) 与 \(c\) 同色''分别属于其中一边。

\subsection{上界靠涂，下界靠找}

最省事的构造是贪心：给顶点排个序，轮到谁就涂上邻居没用过的最小编号颜色。某个顶点最多有 \(\Delta\) 个邻居（\(\Delta\) 是图的最大度），最坏情况下它们占满了 \(\Delta\) 种颜色，第 \(\Delta+1\) 种一定空着。于是

\[
\chi(G)\le\Delta+1 .
\]

这个上界有两处不牢靠。一是结果随顺序漂移：同一个图，按度数从大到小排可能四色搞定，按变量名字母序排可能要六色。贪心不保证顺序好坏，它只保证不越界。二是界本身可能松得离谱。一个中心连着一圈叶子的星形图，\(\Delta\) 等于叶子数，可两色足矣——中心一色，叶子全另一色。

下界从另一头来。图里若有 \(k\) 个顶点两两相邻（一个 \(k\)-团），这 \(k\) 个必须颜色各异，所以 \(\chi(G)\ge\omega(G)\)，\(\omega\) 是最大团的大小。前面那张冲突图正是靠一个三角形锁住了 3。

问题是团并不总能给出真相。

\begin{center}
\begin{tikzpicture}[x=1cm,y=1cm,font=\small,>=stealth]
  \draw[black!60] (0,1.3) -- (-1.236,0.402);
  \draw[black!60] (-1.236,0.402) -- (-0.764,-1.052);
  \draw[black!60] (-0.764,-1.052) -- (0.764,-1.052);
  \draw[red!75,thick] (0.764,-1.052) -- (1.236,0.402);
  \draw[red!75,thick] (1.236,0.402) -- (0,1.3);
  \draw[black!60,fill=blue!30] (0,1.3) circle (7pt);
  \draw[black!60,fill=orange!40] (-1.236,0.402) circle (7pt);
  \draw[black!60,fill=blue!30] (-0.764,-1.052) circle (7pt);
  \draw[black!60,fill=orange!40] (0.764,-1.052) circle (7pt);
  \draw[black!60,fill=white] (1.236,0.402) circle (7pt);
  \node at (1.236,0.402) {\footnotesize ?};
  \node[right,black!70,font=\footnotesize,align=left] at (2.1,0.4)
    {两个邻居一蓝一橙，\\第五个顶点无色可用};
  \node[black!70,font=\footnotesize] at (0,-1.9) {\(C_5\)：无三角形，却要三色};
\end{tikzpicture}

\emph{绕着五边形交替涂蓝、橙，走到最后一个顶点时，它的两个邻居颜色刚好占齐。奇数长度的环合不拢，这是两色着色唯一会失败的方式——也正是上一节判定二分图时那条``没有奇环''的判据，反过来读。}
\end{center}

于是 \(\chi(C_5)=3\) 而 \(\omega(C_5)=2\)。缺口不止这一点点：存在完全不含三角形、色数却要多高有多高的图。障碍未必以``所有人都互相冲突''的形式出现，一条绕不回来的奇环、一圈更隐蔽的结构，同样在消耗颜色。所以下界的工作不能只找团。

实用上，夹逼是常态：找几个团和奇环把下界顶上去，跑几种顶点顺序把上界压下来，两边碰上就收工，碰不上就接受一个区间。精确算 \(\chi\) 是 NP 困难问题，寄存器分配器也不去算它，它们用的是启发式的顶点消去顺序，遇到实在涂不下的变量就把它溢出到内存——把一个顶点从图里删掉，比多要一个寄存器便宜。

顺带记一个跳变：判断 \(\chi\le2\) 只需一遍搜索（上一节的分层染色），判断 \(\chi\le3\) 已经是 NP 完全的。难度不是随颜色数平滑上升的，它在 2 和 3 之间一步跨过去。

\subsection{把冲突画在边上}

换个场景：不是给课程排时段，而是给监考老师排班——同一个老师不能同时出现在两场考试。这时冲突发生在边与边之间：共享端点的两条边不能同色。

这仍然是着色，只是换了个图。把原图 \(G\) 的每条边变成新图的一个顶点，共享端点的两条边在新图里连边，得到线图 \(L(G)\)，于是 \(G\) 的边色数就是 \(L(G)\) 的点色数。语言统一了，但性质不会跟着搬家：Vizing 定理说简单图的边色数只可能是 \(\Delta\) 或 \(\Delta+1\)，二者必居其一。点着色没有这种紧到只剩两个候选的结论。转换省下的是概念，不是信息。

平面图那一头也有类似的落差。四色定理保证任何平面图四色足够，证明极其繁重，是第一个大规模借助计算机完成的主要定理。但它的前提硬得很：图得画得出来且边不交叉。排课冲突图几乎从不平面——只要有五门课两两共享学生，\(K_5\) 就出现了，四色的保证当场失效。定理的力量来自它对结构的强限制，而现实里的冲突图通常不具备那种结构。至于一张图什么时候才真的画得下，留给本章的\hyperref[sec:04-Discrete-Mathematics/03-Graph-Theory/09]{``平面图与 Euler 公式''}。

\subsection{模型只编码了一句话}

图着色能表达的约束只有一条：这两个不能相同。凡是能翻译成``再加一条边''或者``这个顶点的可选颜色集合受限''的要求，框架都装得下——某门课不能排在周一，就相当于给那个顶点禁用一种颜色，这叫列表着色。

越过这条线的约束就装不下了。考试时长不等、教室有容量、老师偏好连排、任务之间有先后顺序、总成本要最小——这些涉及数值、次序或全局目标的要求，图里没有位置放。此时着色仍然是问题的骨架（冲突关系还在），但求解要换成整数规划或约束满足。建模的判断力有一半在这里：认出自己的问题在着色框架内能走多远，以及在哪一步必须换工具。

还有一层假设更隐蔽。着色和上一节的匹配都默认存在一个中央调度者，他掌握全部冲突信息，并且所有人服从安排。如果被安排的是有主意的人——学生挑导师、医学生挑医院、用户挑平台——目标就不再是``总量最大''或``颜色最少''，而是``没有人有动机私下另找搭档''。下一节换成这套设定。
